% Part: sets-functions-relations
% Chapter: sets
% Section: subsets

\documentclass[../../../include/open-logic-section]{subfiles}

\begin{document}

\olfileid{sfr}{set}{sub}
\olsection{子集与幂集}

\begin{explain}
我们经常需要比较集合。一种显而易见的比较方式是：\emph{一个集合中的所有元素也都是另一个集合的元素}。这种情况十分重要，值得引入一些新记号。
\end{explain}

\begin{defn}[子集]
如果集合 $A$ 的每个元素都是 $B$ 的元素，那么我们称 $A$ 是 $B$ 的\emph{子集}，记作 $A \subseteq B$。如果 $A$ 不是 $B$ 的子集，记作 $A \not\subseteq B$。
如果 $A \subseteq B$ 但 $A \neq B$，我们记作 $A \subsetneq B$，并称 $A$ 是 $B$ 的\emph{真子集}。
\end{defn}

\begin{ex}
每个集合都是它自身的子集，且 $\emptyset$ 是任何集合的子集。自然偶数集是自然数集的子集。此外，$\{ a, b \} \subseteq \{ a, b, c \}$。但是 $\{ a, b, e \}$ 不是 $\{ a, b, c \}$ 的子集。
\end{ex}

\begin{ex}
数 $2$ 是整数集的\emph{元素}，而偶数集则是整数集的\emph{子集}。然而，一个集合可能\emph{同时}是另一个集合的元素和子集，例如 $\{0\} \in \{0, \{0\}\}$，以及 $\{0\} \subseteq \{0,
\{0\}\}$。
\end{ex}

外延性给出了集合的同一性标准：$A = B$ 当且仅当 $A$ 的每个元素也是 $B$ 的元素，反之亦然。
“子集”的定义恰好将此标准的前半部分作为 $A \subseteq B$ 的定义：$A$ 的每个元素也是 $B$ 的元素。
当然，如果我们交换 $A$ 和 $B$，该定义也适用：即 $B \subseteq A$ 当且仅当 $B$ 的每个元素也是 $A$ 的元素。而这恰好是外延性中“反之亦然”的部分。换言之，外延性意味着：两个集合相等当且仅当它们互为子集。

\begin{prop}
$A = B$ 当且仅当 $A \subseteq B$ 且 $B \subseteq A$。
\end{prop}

现在也是引入一些有用的辅助记号的好时机。在定义 $A$ 为 $B$ 的子集时，我们说“$A$ 的每个元素都是……”，并在“……”处填上了“$B$ 的一个元素”。但这种表达模式非常常见，为此引入形式记号将大有帮助。

\begin{defn}\ollabel{forallxina}
$(\forall x \in A)\phi$ 是 $\forall x(x \in A \lif \phi)$ 的缩写。类似地，$(\exists x \in A)\phi$ 是 $\exists x(x \in A \land \phi)$ 的缩写。
\end{defn}

使用这个记号，我们可以说 $A \subseteq B$ 当且仅当 $(\forall x \in A)x \in B$。

现在我们来考虑一种特定的集合：一个给定集合的所有子集构成的集合。

\begin{defn}[幂集]
由一个集合~$A$的所有子集构成的集合，称为~$A$的\emph{幂集}，记作 $\Pow{A}$。
  \[
    \Pow{A} = \Setabs{B}{B \subseteq A} 
  \]
\end{defn}

\begin{ex}
$\{ a, b, c \}$ 的所有可能子集有哪些？它们是：
$\emptyset$, $\{a \}$, $\{b\}$, $\{c\}$, $\{a, b\}$, $\{a, c\}$, $\{b,
c\}$, $\{a, b, c\}$。所有这些子集构成的集合是
$\Pow{\{a,b,c\}}$：
\[
\Pow{\{ a, b, c \}} = \{\emptyset, \{a \}, \{b\}, \{c\}, \{a, b\},
\{b, c\}, \{a, c\}, \{a, b, c\}\}
\]
\end{ex}

\begin{prob}
列出 $\{a, b, c, d\}$ 的所有子集。
\end{prob}

\begin{prob}
证明：若 $A$ 有 $n$ 个元素，则 $\Pow{A}$ 有 $2^n$ 个元素。
\end{prob}

\end{document}